% =====================================================================
% TruthfulRAG 模块（Sean：方法分析；Hakimi：实验分析）
% 主入口：thesis.tex 中 % =====================================================================
% TruthfulRAG 模块（Sean：方法分析；Hakimi：实验分析）
% 主入口：thesis.tex 中 % =====================================================================
% TruthfulRAG 模块（Sean：方法分析；Hakimi：实验分析）
% 主入口：thesis.tex 中 % =====================================================================
% TruthfulRAG 模块（Sean：方法分析；Hakimi：实验分析）
% 主入口：thesis.tex 中 \include{body/cl/main}
% 每个小节一个文件，文件头注明负责人；只改自己负责的文件。
% =====================================================================
\providecommand{\cltodo}[1]{\textcolor{red}{[待补：#1]}}
% 图片若尚未生成，则显示占位框，保证随时可编译
\providecommand{\clfig}[2]{\IfFileExists{#1}{\includegraphics[width=#2]{#1}}{\fbox{\parbox{#2}{\centering\cltodo{图片 \texttt{\detokenize{#1}} 尚未生成}}}}}

\section{TruthfulRAG：基于知识图谱的 RAG 事实级冲突消解}
\label{sec:cl}

\input{body/cl/s1_overview}               % X.1   Sean
\input{body/cl/s2_method}                 % X.2   Sean
\input{body/cl/s3_reproduction}           % X.3   Hakimi

\subsection{探究实验}
\label{sec:cl-probe}
本节的实验围绕第~\ref{sec:cl-method-summary}~小节提出的四个假设展开：
\ref{sec:cl-e1}、\ref{sec:cl-e3}~小节检验方法自身的设计假设，
\ref{sec:cl-r3}、\ref{sec:cl-r4}~小节检验论文实验结论的可靠性与可推广性。
\input{body/cl/s4a_probe_method}          % X.4.1–X.4.2  Sean
\input{body/cl/s4b_probe_experiment}      % X.4.3–X.4.4  Hakimi

\subsection{综合讨论：局限与改进方向}
\label{sec:cl-discussion}
\input{body/cl/s5a_discussion_method}     % X.5 前半  Sean
\input{body/cl/s5b_discussion_experiment} % X.5 后半  Hakimi

\input{body/cl/s6_summary}                % X.6   Hakimi

% 每个小节一个文件，文件头注明负责人；只改自己负责的文件。
% =====================================================================
\providecommand{\cltodo}[1]{\textcolor{red}{[待补：#1]}}
% 图片若尚未生成，则显示占位框，保证随时可编译
\providecommand{\clfig}[2]{\IfFileExists{#1}{\includegraphics[width=#2]{#1}}{\fbox{\parbox{#2}{\centering\cltodo{图片 \texttt{\detokenize{#1}} 尚未生成}}}}}

\section{TruthfulRAG：基于知识图谱的 RAG 事实级冲突消解}
\label{sec:cl}

% 论文概述 —— 负责人：Sean
\subsection{论文概述}
\label{sec:cl-overview}

本节依次概述论文的研究背景、问题定义、方法框架以及实验设置与结论，内容均依据论文原文整理。
本组对方法的分析与评价集中在第~\ref{sec:cl-method}~节。

% ---------------------------------------------------------------------
\subsubsection{研究背景}

检索增强生成（Retrieval-Augmented Generation, RAG）将外部检索得到的文档拼接进提示，用以弥补大语言模型参数知识过时或缺失的不足\upcite{lewis2020retrieval,gao2023retrieval}。
随着外部知识库不断更新，检索内容与模型在预训练阶段记住的知识可能互相矛盾，这一现象称为知识冲突。
已有研究表明，发生冲突时模型往往倾向于坚持自身的记忆而忽视检索证据\upcite{longpre2021entity,xie2023adaptive}，从而降低 RAG 系统输出的准确性。

论文将已有的冲突消解方法归为两类。
词元级方法在解码阶段调整输出词元的概率分布，例如对比有无上下文时的输出分布\upcite{yuan2024discerning,jin2024tug}；
语义级方法在文段或陈述层面对齐内部与外部知识，例如 FaithfulRAG 先让模型显式写出自身的参数知识，再与检索内容比对\upcite{zhang2025faithfulrag}。
作者认为两类方法依赖的表示较为零散，难以准确刻画事实之间的关联，因而无法精确定位发生冲突的具体事实。

% ---------------------------------------------------------------------
\subsubsection{问题定义}

记 $\mathcal{M}$ 为具有参数知识 $\mathcal{K}_p$ 的大语言模型。
给定问题 $q$，RAG 系统从外部知识库 $\mathcal{K}_e$ 中检索文档，并以问题和检索内容为条件生成回答。
若存在与 $q$ 相关的两条事实陈述 $f_p\in\mathcal{K}_p$ 与 $f_e\in\mathcal{K}_e$，且二者在事实上不一致，则称发生了知识冲突。
论文采取外部知识优先的立场，目标是在冲突发生时使模型遵循检索内容作答，同时在无冲突时不损害原有性能。
围绕这一目标，论文提出了三个技术问题：如何表示事实知识以便检测冲突，如何为给定问题检索并排序相关事实，以及如何使模型在面对冲突证据时做出可靠决策。

% ---------------------------------------------------------------------
\subsubsection{方法框架}

TruthfulRAG 由图构建、图检索和冲突消解三个模块串联而成，整体流程如图~\ref{fig:cl-pipeline}~所示。

\begin{figure}[tbp]
  \centering
  \clfig{figure/cl_method/pipeline.png}{0.92\textwidth}
  \caption{TruthfulRAG 的整体流程，依次为图构建、图检索与基于熵的冲突消解（取自论文开源仓库）}
  \label{fig:cl-pipeline}
\end{figure}

图构建模块将检索内容划分为语义连贯的文段，由模型 $\mathcal{M}$ 从每个文段中抽取三元组 $(h,r,t)$，汇总得到知识图谱 $\mathcal{G}=(\mathcal{E},\mathcal{R},\mathcal{T})$，其中 $\mathcal{E}$、$\mathcal{R}$、$\mathcal{T}$ 分别为实体、关系与三元组集合。

图检索模块先从问题中抽取关键要素，按语义相似度选出最相关的实体集合 $\mathcal{E}_{imp}$ 与关系集合 $\mathcal{R}_{imp}$，再从每个关键实体出发做两跳遍历，得到初始路径集合 $\mathcal{P}_{init}$。
每条路径按其对关键实体与关系的覆盖程度打分：
\begin{equation}
  \mathrm{Ref}(p)=\alpha\,\frac{|\{e\in p\}\cap\mathcal{E}_{imp}|}{|\mathcal{E}_{imp}|}
               +\beta\,\frac{|\{r\in p\}\cap\mathcal{R}_{imp}|}{|\mathcal{R}_{imp}|},
  \label{eq:cl-ref}
\end{equation}
得分最高的 $K$ 条路径构成 $\mathcal{P}_{super}$。每条路径被展开为路径序列、实体属性与关系属性三部分文本。

冲突消解模块以输出熵度量模型的不确定性。设生成答案长度为 $L$，第 $t$ 个位置上前 $k$ 个候选词元的概率为 $p_i^{(t)}$，则
\begin{equation}
  H = -\frac{1}{L}\sum_{t=1}^{L}\sum_{i=1}^{k} p_i^{(t)}\log_2 p_i^{(t)} .
  \label{eq:cl-entropy}
\end{equation}
分别计算只给出问题时的熵 $H_{param}$ 与给出问题和路径 $p$ 时的熵 $H_{aug}(p)$，二者之差 $\Delta H_p=H_{aug}(p)-H_{param}$ 被视为冲突的指示量。
论文认为 $\Delta H_p$ 为正说明该路径与模型的参数知识相矛盾，对模型具有纠正作用，因此保留
\begin{equation}
  \mathcal{P}_{corrective}=\{\,p\in\mathcal{P}_{super}\mid \Delta H_p>\tau\,\},
  \label{eq:cl-paperrule}
\end{equation}
并以 $\mathcal{M}(q\oplus\mathcal{P}_{corrective})$ 生成最终回答，其中 $\tau$ 为阈值，$\oplus$ 表示文本拼接。

% ---------------------------------------------------------------------
\subsubsection{实验设置与主要结论}

论文在四个数据集上评测：FaithEval\upcite{ming2025faitheval} 包含反事实与不一致的上下文；MuSiQue\upcite{trivedi2022musique} 与 SQuAD\upcite{rajpurkar2016squad} 采用 KRE\upcite{ying2024intuitive} 构造的反事实版本；RealtimeQA\upcite{kasai2023realtime} 侧重时效性冲突。
所用模型为 GPT-4o-mini、Qwen2.5-7B-Instruct 与 Mistral-7B-Instruct。
对比基线包括不使用检索的直接生成、标准 RAG、提示优化方法 KRE、解码干预方法 COIECD\upcite{yuan2024discerning} 以及 FaithfulRAG，评价指标为准确率。

论文报告的主要结论有四点。
其一，TruthfulRAG 在三个模型上均取得最高的平均准确率，在 GPT-4o-mini、Qwen2.5-7B 与 Mistral-7B 上分别为 78.8\%、78.3\% 与 81.3\%，次优的 FaithfulRAG 分别为 76.5\%、78.1\% 与 80.7\%。
其二，在 MuSiQue 与 SQuAD 的无冲突版本上，该方法的准确率不低于标准 RAG。
其三，与自然语言上下文相比，结构化的推理路径使模型赋予正确答案更高的对数概率。
其四，消融实验表明图模块与冲突消解模块均对性能有贡献。
               % X.1   Sean
% X.2 方法剖析 —— 负责人：Sean
\subsection{方法剖析}
\label{sec:cl-method}

本节的分析同时依据论文正文和作者开源的代码。我们逐函数阅读了实现，下文提到的代码位置均以“文件名中的函数名”标注，便于核对。
分析按“问题形式化 → 三个模块 → 论文与实现的对照 → 优缺点与待检验假设”的顺序展开。

% ---------------------------------------------------------------------
\subsubsection{问题形式化与设计动机}
\label{sec:cl-formal}

记 LLM 为 $\mathcal{M}$，问题为 $q$，检索到的上下文为 $C$。
模型只凭参数知识作答时的输出分布为 $p_{\mathcal{M}}(y\mid q)$，给定上下文后为 $p_{\mathcal{M}}(y\mid q, C)$。
令 $a_{\text{param}}=\arg\max_y p_{\mathcal{M}}(y\mid q)$ 为模型“记住的”答案，$a_{\text{ctx}}$ 为上下文 $C$ 支持的答案。
\textbf{知识冲突}即 $a_{\text{param}}\neq a_{\text{ctx}}$ 的情形。论文采取“外部知识优先”的立场，目标是让最终输出等于 $a_{\text{ctx}}$。
这一立场隐含了一个前提：外部知识是对的。我们在下文把它记作假设~A1。

\paragraph{为什么是“事实级”。}
我们对这一动机的理解是：冲突往往只集中在\textbf{极少数}事实上。
以论文附录的例子为例，检索文本把“新拉雷多市所在的州”从 Tamaulipas 改成了 Sinaloa，而同一段文字中其余几十条事实都与模型的记忆一致。
在整段文本上做“语义级”比较时，这一条冲突会被大量一致的信息稀释；在输出词元上做“词元级”调整时，方法又并不知道是哪条事实引起了分歧。
若把文本拆成三元组，冲突就被定位到一条边（Nuevo Laredo, located\_in, Sinaloa）上，可以单独处理。

\paragraph{方法的实质。}
值得注意的是，TruthfulRAG 并\textbf{没有}显式地把模型的参数知识写成三元组、再与外部三元组逐条比对（FaithfulRAG 正是这样做的\upcite{zhang2025faithfulrag}）。
它采用的是一个\textbf{间接探针}：如果把某条路径加入提示后模型变得更不确定，就推断这条路径与模型的信念相矛盾。
因此，我们把这个方法概括为三步：“\textbf{结构化重写} + \textbf{基于不确定性的路径选择} + \textbf{提示增强}”。
它能否奏效，取决于以下四个隐含假设：
\begin{description}
  \item[A1 外部知识可信。] 冲突时外部知识总是对的。若检索内容本身有误（如被投毒或已过时），方法会主动放大错误。
  \item[A2 熵升高意味着“纠正性”。] 加入路径后熵升高 $\Leftrightarrow$ 路径与参数知识冲突 $\Leftrightarrow$ 路径对答题有用。但无关或嘈杂的信息同样会让模型变得不确定；反之，一条与模型记忆\emph{一致}的正确路径会降低熵，却同样是答题所需的。
  \item[A3 三元组抽取是忠实的。] 抽取由同一个 LLM 完成。面对反事实文本，模型可能用自己的先验“悄悄改正”原文，或干脆漏抽，这样冲突在抽取阶段就已经被抹掉了。
  \item[A4 答案位于查询实体的两跳之内。] 检索只做两跳遍历，而 MuSiQue 等多跳问题可能需要三到四跳。
\end{description}

% ---------------------------------------------------------------------
\subsubsection{三个模块的逐一分析}
\label{sec:cl-modules}

\paragraph{（1）图构建。}
论文的描述是：先对检索内容做语义分段，再用 LLM 逐段抽取三元组 $(h,r,t)$，汇总成图 $\mathcal{G}=(\mathcal{E},\mathcal{R},\mathcal{T})$。
代码中的实现（\texttt{modules.py} 中的 \texttt{generate\_knowledge\_graph}）是：
按 1024 个 token、重叠 128 的窗口切分文本；
对每个分块调用一次抽取提示，再追加一轮“继续抽取”（gleaning），即每个分块固定 2 次 LLM 调用；
输出的实体带有名称、类型和一段自由文本描述，关系带有描述、关键词和权重；
同名实体合并，描述拼接。
这一流程与 LightRAG\upcite{guo2024lightrag} 的图构建实现高度一致。
我们注意到三点：
\begin{itemize}
  \item \textbf{图是“每题局部”的}：每道题都只用这道题自己的检索文本从零建一张图，答完即清空（\texttt{pipeline.py} 中的 \texttt{make\_knowledge\_graph}）。因此它并非一个持久的知识库，建图成本无法在多道题之间摊销，也无法利用其它文档做交叉验证。
  \item \textbf{实体对齐只靠名称完全匹配}（统一大写），别名、指代不会合并，同一实体可能分裂成多个节点。
  \item \textbf{“结构化”是有限度的}：节点和边上挂的都是自然语言描述。最终交给模型的路径文本（见下）大部分仍是自然语言句子，只是按实体和关系组织了起来。
\end{itemize}

\paragraph{（2）图检索。}
论文用式（\ref{eq:cl-ref}）对候选路径打分，并取前 $K$ 条：
\begin{equation}
  \mathrm{Ref}(p)=\alpha\cdot\frac{|\{e\in p\}\cap\mathcal{E}_{imp}|}{|\mathcal{E}_{imp}|}
               +\beta\cdot\frac{|\{r\in p\}\cap\mathcal{R}_{imp}|}{|\mathcal{R}_{imp}|}
  \label{eq:cl-ref}
\end{equation}
代码（\texttt{modules.py} 中的 \texttt{\_build\_query\_context}）的实际流程是：
\begin{enumerate}
  \item 用 LLM 从问题中抽取至多 5 个实体，以及“答案类型”关键词（从图中已有的实体类型里挑选）；
  \item 对每个问题实体，在实体名称向量库中检索最相似的 10 个节点（余弦相似度需大于 0.2），作为起点；
  \item 从每个起点枚举所有两跳路径；找不到路径的孤立起点只保留相似度最高的 20\%；
  \item 统计每条路径上属于“答案类型”的节点个数；
  \item 用整个问题在关系向量库中检索若干条边，其中与关键实体相连的边记为“有效边”；一条路径只要包含至少一条有效边，就被保留；
  \item 删除被其它路径包含的子路径，再把每条路径渲染为“Path / Nodes / Edges”三段文本。
\end{enumerate}
这套逻辑与 MiniRAG\upcite{fan2025minirag} 的异构图检索在函数结构和提示模板上高度一致。
关键在于：第 4、5 步算出的计数\textbf{只用于判断去留，并不用于排序}。最终路径列表按路径长度排列，随后由 \texttt{pipeline.py} 中的 \texttt{run} 直接截取前 30 条送入下一模块。
换言之，式（\ref{eq:cl-ref}）中的 $\alpha$、$\beta$ 和按得分取前 $K$ 条的操作在代码中并不存在，“选哪几条路径”实际上完全由下一步的熵过滤决定。

\paragraph{（3）基于熵的冲突消解。}
论文用式（\ref{eq:cl-entropy}）衡量模型输出的不确定性：
\begin{equation}
  H\big(P(\cdot\mid \text{ctx})\big) = -\frac{1}{L}\sum_{t=1}^{L}\sum_{i=1}^{k} p_i^{(t)}\log p_i^{(t)},
  \qquad
  \Delta H_p = H\big(P(\cdot\mid q,p)\big) - H\big(P(\cdot\mid q)\big),
  \label{eq:cl-entropy}
\end{equation}
并保留 $\Delta H_p>\tau$ 的路径。代码实现（\texttt{modules.py} 中的 \texttt{entropy\_filter}，\texttt{llm/backend.py} 中的 \texttt{cal\_entroy}）有以下细节：
\begin{itemize}
  \item 计算 $H(\cdot\mid q)$ 和 $H(\cdot\mid q,p)$ 时，提示中\textbf{不含原始上下文}，只有问题（和选项）以及该条路径；用贪心解码生成，取每个位置概率最高的 10 个词元，重新归一化后按\textbf{自然对数}计算熵（论文写的是 $\log_2$）。
  \item 生成长度 $L$ 被设为\textbf{金标准答案的 token 数}。也就是说，推理阶段用到了标注信息，存在轻微的标签泄漏。熵是在模型自己生成的前 $L$ 个词元上计算的，而不是在正确答案上。
  \item 每条候选路径都要单独调用一次模型，每道题共需 $1+|P|$ 次调用（$|P|\le 30$）。
\end{itemize}
最关键的差异在于选择规则。代码实际执行的是：
\begin{equation}
  \text{keep}(p) = \big[\Delta H_p \ge 0\big] \;\lor\; \Big[\,H_{aug}(p)>0 \;\land\; \frac{H_{base}}{H_{aug}(p)} \le 10^{\tau}\Big],
  \label{eq:cl-coderule}
\end{equation}
然后按 $\Delta H_p$ 从大到小取前 10 条；若一条都不满足，则保留熵最大的那条。
按式（\ref{eq:cl-coderule}），当 $\tau=1$ 时，一条使熵\emph{下降}的路径只有在熵降到原来的十分之一以下时才会被丢弃；当 $\tau=3$ 时，则要降到千分之一以下。
因此这条规则几乎等价于“\textbf{全部保留，按 $\Delta H$ 排序取前 10}”，$\tau$ 的作用非常弱。
这也为论文附录中“$\tau=1$ 与 $\tau=3$ 结果几乎相同”的现象提供了另一种解释：它未必说明方法对超参数鲁棒，更可能只是 $\tau$ 本身几乎不起作用。
反过来，如果真的按论文写的 $\Delta H_p>\tau$ 执行：在 top-10 重归一化下，每个位置的熵上限为 $\ln 10\approx 2.30$，$\tau=1$ 已是很严格的要求，很多题可能一条路径都留不下来（第~\ref{sec:cl-e1}~小节将验证这一点）。

\paragraph{（4）最终生成。}
论文写作 $\text{Response}=\mathcal{M}(q\oplus\mathcal{P}_{corrective})$，读者容易理解为用筛选后的路径\emph{替换}原始上下文。
代码（\texttt{modules.py} 中的 \texttt{predict\_answer}）实际构造的提示同时包含筛选后的路径（标注为 Facts）\textbf{和完整的原始上下文}（标注为 Context），并要求模型“优先依据 Facts，信息不足时再参考 Context”，以思维链加 JSON 的格式输出。
所以更准确的写法是 $\mathcal{M}(q\oplus\mathcal{P}_{kept}\oplus C)$：方法是在原始上下文之上\textbf{追加}一组被强调的事实，而不是替换上下文。
这一点直接影响对论文消融实验和“上下文长度”统计的解读。

% ---------------------------------------------------------------------
\subsubsection{论文描述与开源实现的对照}
\label{sec:cl-gap}

表~\ref{tab:cl-gap}~汇总了上文发现的差异。需要说明的是，开源代码未必就是论文实验所用的版本，因此我们不把这些差异直接当作论文的错误，而是将它们作为后续实验需要控制或检验的因素。

\begin{table}[htbp]
  \centering
  \small
  \caption{论文描述与开源实现的主要差异}
  \label{tab:cl-gap}
  \begin{tabular}{p{1.6cm}p{3.6cm}p{4.4cm}p{4.2cm}}
    \toprule
    环节 & 论文描述 & 代码实现 & 影响 \\
    \midrule
    路径打分 & 式（\ref{eq:cl-ref}），按 $\mathrm{Ref}(p)$ 取前 $K$ 条 & 无 $\alpha,\beta$；计数只决定去留，按路径长度排序后截取前 30 条 & 检索阶段不做相关性排序，选择完全依赖熵过滤 \\
    熵的计算 & $\log_2$，答案的 token 分布 & top-10 重归一化，自然对数；生成长度 = 金标准答案长度 & 推理时使用了标注信息；阈值的量纲与论文不同 \\
    过滤规则 & 保留 $\Delta H_p>\tau$ & 式（\ref{eq:cl-coderule}），再取前 10 条；为空时保留熵最大者 & 几乎全部保留，$\tau$ 作用很弱 \\
    最终输入 & $q\oplus\mathcal{P}_{corrective}$ & $q\oplus\mathcal{P}_{kept}\oplus C$，提示要求优先依据 Facts & 是“增强”而非“替换” \\
    知识图谱 & 知识图谱 & 每题从零构建的局部图 & 建图成本无法摊销，没有跨文档校验 \\
    评测指标 & 准确率 ACC & 归一化后的\emph{预测}是\emph{金标准}的子串即判对；空预测也判对 & 指标偏宽松；我们同时报告严格的精确匹配（EM） \\
    \bottomrule
  \end{tabular}
\end{table}

% ---------------------------------------------------------------------
\subsubsection{方法层面的优缺点小结}
\label{sec:cl-method-summary}

\paragraph{优点。}
（1）\textbf{定位粒度合适}：三元组把冲突限定在一条边上，路径文本比原文更短、噪声更少，这与论文图 3 中模型对正确答案更自信的观察相符。
（2）\textbf{与模型无关、即插即用}：只要能拿到输出词元的对数概率，就能接入任意 RAG 系统，不需要训练。
（3）\textbf{中间结果可审计}：每条候选路径和它的 $\Delta H$ 都可以记录下来，便于定位错误发生在哪一步。我们正是利用这一点设计了逐级漏斗分析（第~\ref{sec:cl-e3}~小节）。

\paragraph{缺点。}
（1）\textbf{成本高}：每道题需要“$2\times$ 分块数（建图）$+1$（关键词）$+1+|P|$（熵）$+1$（生成）”次调用。论文附录给出的单题耗时为 36--62 秒，而标准 RAG 仅需 0.4--2.6 秒。
（2）\textbf{核心信号缺乏直接验证}：A2 是整个方法的立足点，论文只给出了最终准确率，没有直接检验“熵升高的路径是否真的包含正确答案”。
（3）\textbf{对错误外部知识没有防御}：在 A1 不成立时，方法会把错误信息以“Facts”的形式加以强调。
（4）\textbf{误差逐级传递}：抽取漏掉或“改正”的事实，后续模块无从恢复（A3、A4）。
（5）\textbf{依赖标注与对数概率}：生成长度用到了金标准答案；而很多闭源 API 并不提供对数概率。

\paragraph{待检验的假设。}
基于以上分析，我们提出四个可检验的假设，并与后续实验一一对应，如表~\ref{tab:cl-hypo}~所示。

\begin{table}[htbp]
  \centering
  \small
  \caption{由方法分析得到的假设与对应实验}
  \label{tab:cl-hypo}
  \begin{tabular}{p{0.6cm}p{8.6cm}p{2.0cm}p{1.6cm}}
    \toprule
    编号 & 假设 & 对应前提 & 实验 \\
    \midrule
    H1 & $\Delta H$ 区分“含答案路径”与“不含答案路径”的能力有限 & A2 & E1 \\
    H2 & 按论文规则过滤会丢掉大量有用路径；代码规则的收益主要来自排序与提示增强，而非“过滤冲突” & A2、实现差异 & E2 \\
    H3 & 答案主要在抽取、检索阶段丢失，而非在过滤阶段 & A3、A4 & E3 \\
    H4 & 当外部知识本身有误时，方法比标准 RAG 更容易盲从 & A1 & R3 \\
    \bottomrule
  \end{tabular}
\end{table}
                 % X.2   Sean
% 复现实验 —— 负责人：Hakimi
% （占位文件：由 Hakimi 撰写。标签 sec:cl-repro 已被其它小节引用，请保留。）
\subsection{复现实验}
\label{sec:cl-repro}

\subsubsection{实验设置}
\cltodo{Hakimi：子集抽样方式（seed=42）、模型、运行环境、对代码所做的改动（逐题记录、并发）、成本}

\subsubsection{结果与对比}
\cltodo{Hakimi：R1、R2 结果，并与论文表 1、表 2 对比}

\subsubsection{原文实验设计评述}
\cltodo{Hakimi：显著性检验方式、CPR 定义、Qwen-72B 上的异常值、RealtimeQA 样本量}
           % X.3   Hakimi

\subsection{探究实验}
\label{sec:cl-probe}
本节的实验围绕第~\ref{sec:cl-method-summary}~小节提出的四个假设展开：
\ref{sec:cl-e1}、\ref{sec:cl-e3}~小节检验方法自身的设计假设，
\ref{sec:cl-r3}、\ref{sec:cl-r4}~小节检验论文实验结论的可靠性与可推广性。
% X.4.1–X.4.2 探究实验（方法假设）—— 负责人：Sean
% 数据来源：results/trace_*.jsonl（Hakimi 运行 R1/R2 时产生）
% 脚本：experiments/analysis/e1_entropy.py, e2_filters.py, e3_funnel.py
% 图：figure/cl_method/ ；表：experiments/out/cl_method/

\paragraph{共同设置。}
本节实验全部复用复现实验（第~\ref{sec:cl-repro}~小节）中 TruthfulRAG 的运行记录。
我们在开源代码中加入了逐题记录，保存每道题的知识图谱规模、全部候选路径及其熵值 $H$ 与 $\Delta H$、被保留的路径编号以及最终答案。
记录逻辑只在原有计算之外附加输出，不改变方法的行为。我们还按式（\ref{eq:cl-coderule}）独立重写了一遍选择规则，在离线测试中，它选出的路径与原实现记录的路径完全一致；E2 中的“代码规则”正是基于这一重写版本。
一条路径“含答案”的判定方法是：按官方评测脚本的规则归一化后，金标准答案作为完整词组出现在路径文本中。
对反事实数据集（MuSiQue、SQuAD），金标准是检索文本中被篡改后的答案，也就是方法应当“跟随”的那个答案。
模型为 GPT-4o-mini，每个数据集抽取样本 \cltodo{N}~条。

% ---------------------------------------------------------------------
\subsubsection{熵过滤机制是否成立（E1、E2）}
\label{sec:cl-e1}

\paragraph{E1：$\Delta H$ 能否识别“纠正性路径”。}
如果假设 A2 成立，含答案的路径应当系统性地具有更大的 $\Delta H$。
我们把所有候选路径按“是否含答案”分为两组，比较两组 $\Delta H$ 的分布，并报告三个指标：
（1）以 $\Delta H$ 为打分、区分两组的 AUC（0.5 即等同于随机）；
（2）单侧 Mann--Whitney U 检验的 $p$ 值；
（3）在至少有一条含答案路径的题目中，$\Delta H$ 最大的那条路径恰好含答案的比例（Hit@1），并与随机排序下的期望值对比。

\begin{figure}[htbp]
  \centering
  \clfig{figure/cl_method/e1_dh_by_label.pdf}{0.85\textwidth}
  \caption{含答案路径与不含答案路径的 $\Delta H$ 分布（横线为中位数，虚线为 $\Delta H=0$）}
  \label{fig:cl-e1}
\end{figure}

\cltodo{表：各数据集的路径数、含答案比例、两组平均 $\Delta H$、AUC、$p$ 值、Hit@1（来自 \texttt{e1\_entropy.tex}）}

\cltodo{结果分析：AUC 是否显著高于 0.5；含答案路径中 $\Delta H>0$ 的比例；在 golden（无冲突）数据上，含答案路径的 $\Delta H$ 是否转为负值——若是，则说明熵升高反映的是“与记忆冲突”，而不是“对答题有用”，从而支持 H1}

\paragraph{E2：过滤策略对比。}
候选路径及其熵值均已记录，因此更换过滤规则时只需重做最后一步生成（每题 1 次调用）。我们比较以下策略：
\begin{itemize}
  \item \textbf{代码规则}：式（\ref{eq:cl-coderule}），即原始实现；
  \item \textbf{论文规则}：只保留 $\Delta H_p>\tau$ 的路径；
  \item \textbf{全部保留}：不过滤，对应代码自带的消融设置；
  \item \textbf{随机}：随机保留与代码规则同样数量的路径；
  \item \textbf{反向}：保留 $\Delta H$ 最\emph{小}的同样数量的路径；
  \item \textbf{不加路径}：只用问题和原始上下文，提示格式不变。
\end{itemize}
其中，“随机”和“反向”用来检验 $\Delta H$ 排序本身是否有价值；“不加路径”用来衡量整个知识图谱分支的净贡献。
我们还在 $\tau\in\{0,0.5,1,2,3\}$ 上扫描论文规则，统计“保留路径中含答案”和“一条路径都没保留”的题目比例。

\begin{figure}[htbp]
  \centering
  \begin{minipage}{0.48\textwidth}\centering
    \clfig{figure/cl_method/e2_accuracy.pdf}{\linewidth}
  \end{minipage}\hfill
  \begin{minipage}{0.48\textwidth}\centering
    \clfig{figure/cl_method/e2_tau_coverage.pdf}{\linewidth}
  \end{minipage}
  \caption{左：不同过滤策略的平均准确率；右：论文规则下阈值 $\tau$ 对保留结果的影响}
  \label{fig:cl-e2}
\end{figure}

\cltodo{表：各策略在各数据集上的准确率（来自 \texttt{e2\_accuracy.tex}）}

\cltodo{结果分析：代码规则与“随机”、“反向”、“不加路径”之间的差距；论文规则在 $\tau=1$ 时空集的比例；据此判断 H2}

% ---------------------------------------------------------------------
\subsubsection{流水线逐级损失与错误归因（E3）}
\label{sec:cl-e3}

为定位答案在哪一步丢失，我们逐题检查答案是否依次出现在：
S1 原始检索文本、S2 构建出的知识图谱、S3 候选路径、S4 被保留的路径中，以及 S5 最终是否答对。
由于最终提示中仍包含完整的原始上下文（第~\ref{sec:cl-modules}~小节），S5 并不受 S4 约束。
因此我们额外比较“保留路径含答案”与“保留路径不含答案”两种情况下的准确率，以区分模型的正确答案究竟来自路径，还是来自原文。

\begin{figure}[htbp]
  \centering
  \clfig{figure/cl_method/e3_funnel.pdf}{\textwidth}
  \caption{各数据集上答案的逐级保留率（虚线为标准 RAG 的准确率）}
  \label{fig:cl-e3}
\end{figure}

\cltodo{结果分析：最大的一级损失发生在哪里（S1→S2 抽取，还是 S2→S3 检索）；反事实数据上 S1→S2 的损失是否明显高于 golden 数据（即 A3 中的“抽取时被模型先验改正”）；S4 不成立时的准确率}

\cltodo{（P1）人工归因：从答错的样本中抽取 30 条，归入抽取错误、检索遗漏、过滤误删、生成错误四类，并给出一个典型案例}
          % X.4.1–X.4.2  Sean
% 补充实验 R3–R4—— 负责人：Hakimi
% （占位文件：由 Hakimi 撰写。标签 sec:cl-r3 / sec:cl-r4 已被其它小节引用，请保留。）
\subsubsection{外部知识错误下的鲁棒性}
\label{sec:cl-r3}
\cltodo{Hakimi：用 negative 数据的记录、以 gold\_true 重新评估；检验假设 H4（表~\ref{tab:cl-hypo}）}

\subsubsection{统计显著性与计算成本}
\label{sec:cl-r4}
\cltodo{Hakimi：McNemar 检验、bootstrap 置信区间、准确率与成本的关系}
      % X.4.3–X.4.4  Hakimi

\subsection{综合讨论：局限与改进方向}
\label{sec:cl-discussion}
% X.5 前半：方法层面的局限与改进 —— 负责人：Sean（D2 撰写）
\cltodo{Sean：结合 E1–E3 的结果，讨论方法层面的局限，并提出改进方向（例如：用“答案分布是否翻转”替代熵变化作为冲突信号；先在原文中定位冲突句、只抽取相关三元组以降低建图成本；对外部知识增加可信度校验以防御 A1 失效）}
     % X.5 前半  Sean
% 局限与展望：实验层面 —— 负责人：Hakimi
\subsubsection{实验层面}
\label{sec:cl-discussion-exp}
\cltodo{Hakimi：实验层面的局限（样本规模、单一模型、评测指标等）}
 % X.5 后半  Hakimi

% 小结 —— 负责人：Hakimi
\subsection{小结}
\label{sec:cl-summary}
\cltodo{Hakimi：半页小结}
                % X.6   Hakimi

% 每个小节一个文件，文件头注明负责人；只改自己负责的文件。
% =====================================================================
\providecommand{\cltodo}[1]{\textcolor{red}{[待补：#1]}}
% 图片若尚未生成，则显示占位框，保证随时可编译
\providecommand{\clfig}[2]{\IfFileExists{#1}{\includegraphics[width=#2]{#1}}{\fbox{\parbox{#2}{\centering\cltodo{图片 \texttt{\detokenize{#1}} 尚未生成}}}}}

\section{TruthfulRAG：基于知识图谱的 RAG 事实级冲突消解}
\label{sec:cl}

% 论文概述 —— 负责人：Sean
\subsection{论文概述}
\label{sec:cl-overview}

本节依次概述论文的研究背景、问题定义、方法框架以及实验设置与结论，内容均依据论文原文整理。
本组对方法的分析与评价集中在第~\ref{sec:cl-method}~节。

% ---------------------------------------------------------------------
\subsubsection{研究背景}

检索增强生成（Retrieval-Augmented Generation, RAG）将外部检索得到的文档拼接进提示，用以弥补大语言模型参数知识过时或缺失的不足\upcite{lewis2020retrieval,gao2023retrieval}。
随着外部知识库不断更新，检索内容与模型在预训练阶段记住的知识可能互相矛盾，这一现象称为知识冲突。
已有研究表明，发生冲突时模型往往倾向于坚持自身的记忆而忽视检索证据\upcite{longpre2021entity,xie2023adaptive}，从而降低 RAG 系统输出的准确性。

论文将已有的冲突消解方法归为两类。
词元级方法在解码阶段调整输出词元的概率分布，例如对比有无上下文时的输出分布\upcite{yuan2024discerning,jin2024tug}；
语义级方法在文段或陈述层面对齐内部与外部知识，例如 FaithfulRAG 先让模型显式写出自身的参数知识，再与检索内容比对\upcite{zhang2025faithfulrag}。
作者认为两类方法依赖的表示较为零散，难以准确刻画事实之间的关联，因而无法精确定位发生冲突的具体事实。

% ---------------------------------------------------------------------
\subsubsection{问题定义}

记 $\mathcal{M}$ 为具有参数知识 $\mathcal{K}_p$ 的大语言模型。
给定问题 $q$，RAG 系统从外部知识库 $\mathcal{K}_e$ 中检索文档，并以问题和检索内容为条件生成回答。
若存在与 $q$ 相关的两条事实陈述 $f_p\in\mathcal{K}_p$ 与 $f_e\in\mathcal{K}_e$，且二者在事实上不一致，则称发生了知识冲突。
论文采取外部知识优先的立场，目标是在冲突发生时使模型遵循检索内容作答，同时在无冲突时不损害原有性能。
围绕这一目标，论文提出了三个技术问题：如何表示事实知识以便检测冲突，如何为给定问题检索并排序相关事实，以及如何使模型在面对冲突证据时做出可靠决策。

% ---------------------------------------------------------------------
\subsubsection{方法框架}

TruthfulRAG 由图构建、图检索和冲突消解三个模块串联而成，整体流程如图~\ref{fig:cl-pipeline}~所示。

\begin{figure}[tbp]
  \centering
  \clfig{figure/cl_method/pipeline.png}{0.92\textwidth}
  \caption{TruthfulRAG 的整体流程，依次为图构建、图检索与基于熵的冲突消解（取自论文开源仓库）}
  \label{fig:cl-pipeline}
\end{figure}

图构建模块将检索内容划分为语义连贯的文段，由模型 $\mathcal{M}$ 从每个文段中抽取三元组 $(h,r,t)$，汇总得到知识图谱 $\mathcal{G}=(\mathcal{E},\mathcal{R},\mathcal{T})$，其中 $\mathcal{E}$、$\mathcal{R}$、$\mathcal{T}$ 分别为实体、关系与三元组集合。

图检索模块先从问题中抽取关键要素，按语义相似度选出最相关的实体集合 $\mathcal{E}_{imp}$ 与关系集合 $\mathcal{R}_{imp}$，再从每个关键实体出发做两跳遍历，得到初始路径集合 $\mathcal{P}_{init}$。
每条路径按其对关键实体与关系的覆盖程度打分：
\begin{equation}
  \mathrm{Ref}(p)=\alpha\,\frac{|\{e\in p\}\cap\mathcal{E}_{imp}|}{|\mathcal{E}_{imp}|}
               +\beta\,\frac{|\{r\in p\}\cap\mathcal{R}_{imp}|}{|\mathcal{R}_{imp}|},
  \label{eq:cl-ref}
\end{equation}
得分最高的 $K$ 条路径构成 $\mathcal{P}_{super}$。每条路径被展开为路径序列、实体属性与关系属性三部分文本。

冲突消解模块以输出熵度量模型的不确定性。设生成答案长度为 $L$，第 $t$ 个位置上前 $k$ 个候选词元的概率为 $p_i^{(t)}$，则
\begin{equation}
  H = -\frac{1}{L}\sum_{t=1}^{L}\sum_{i=1}^{k} p_i^{(t)}\log_2 p_i^{(t)} .
  \label{eq:cl-entropy}
\end{equation}
分别计算只给出问题时的熵 $H_{param}$ 与给出问题和路径 $p$ 时的熵 $H_{aug}(p)$，二者之差 $\Delta H_p=H_{aug}(p)-H_{param}$ 被视为冲突的指示量。
论文认为 $\Delta H_p$ 为正说明该路径与模型的参数知识相矛盾，对模型具有纠正作用，因此保留
\begin{equation}
  \mathcal{P}_{corrective}=\{\,p\in\mathcal{P}_{super}\mid \Delta H_p>\tau\,\},
  \label{eq:cl-paperrule}
\end{equation}
并以 $\mathcal{M}(q\oplus\mathcal{P}_{corrective})$ 生成最终回答，其中 $\tau$ 为阈值，$\oplus$ 表示文本拼接。

% ---------------------------------------------------------------------
\subsubsection{实验设置与主要结论}

论文在四个数据集上评测：FaithEval\upcite{ming2025faitheval} 包含反事实与不一致的上下文；MuSiQue\upcite{trivedi2022musique} 与 SQuAD\upcite{rajpurkar2016squad} 采用 KRE\upcite{ying2024intuitive} 构造的反事实版本；RealtimeQA\upcite{kasai2023realtime} 侧重时效性冲突。
所用模型为 GPT-4o-mini、Qwen2.5-7B-Instruct 与 Mistral-7B-Instruct。
对比基线包括不使用检索的直接生成、标准 RAG、提示优化方法 KRE、解码干预方法 COIECD\upcite{yuan2024discerning} 以及 FaithfulRAG，评价指标为准确率。

论文报告的主要结论有四点。
其一，TruthfulRAG 在三个模型上均取得最高的平均准确率，在 GPT-4o-mini、Qwen2.5-7B 与 Mistral-7B 上分别为 78.8\%、78.3\% 与 81.3\%，次优的 FaithfulRAG 分别为 76.5\%、78.1\% 与 80.7\%。
其二，在 MuSiQue 与 SQuAD 的无冲突版本上，该方法的准确率不低于标准 RAG。
其三，与自然语言上下文相比，结构化的推理路径使模型赋予正确答案更高的对数概率。
其四，消融实验表明图模块与冲突消解模块均对性能有贡献。
               % X.1   Sean
% X.2 方法剖析 —— 负责人：Sean
\subsection{方法剖析}
\label{sec:cl-method}

本节的分析同时依据论文正文和作者开源的代码。我们逐函数阅读了实现，下文提到的代码位置均以“文件名中的函数名”标注，便于核对。
分析按“问题形式化 → 三个模块 → 论文与实现的对照 → 优缺点与待检验假设”的顺序展开。

% ---------------------------------------------------------------------
\subsubsection{问题形式化与设计动机}
\label{sec:cl-formal}

记 LLM 为 $\mathcal{M}$，问题为 $q$，检索到的上下文为 $C$。
模型只凭参数知识作答时的输出分布为 $p_{\mathcal{M}}(y\mid q)$，给定上下文后为 $p_{\mathcal{M}}(y\mid q, C)$。
令 $a_{\text{param}}=\arg\max_y p_{\mathcal{M}}(y\mid q)$ 为模型“记住的”答案，$a_{\text{ctx}}$ 为上下文 $C$ 支持的答案。
\textbf{知识冲突}即 $a_{\text{param}}\neq a_{\text{ctx}}$ 的情形。论文采取“外部知识优先”的立场，目标是让最终输出等于 $a_{\text{ctx}}$。
这一立场隐含了一个前提：外部知识是对的。我们在下文把它记作假设~A1。

\paragraph{为什么是“事实级”。}
我们对这一动机的理解是：冲突往往只集中在\textbf{极少数}事实上。
以论文附录的例子为例，检索文本把“新拉雷多市所在的州”从 Tamaulipas 改成了 Sinaloa，而同一段文字中其余几十条事实都与模型的记忆一致。
在整段文本上做“语义级”比较时，这一条冲突会被大量一致的信息稀释；在输出词元上做“词元级”调整时，方法又并不知道是哪条事实引起了分歧。
若把文本拆成三元组，冲突就被定位到一条边（Nuevo Laredo, located\_in, Sinaloa）上，可以单独处理。

\paragraph{方法的实质。}
值得注意的是，TruthfulRAG 并\textbf{没有}显式地把模型的参数知识写成三元组、再与外部三元组逐条比对（FaithfulRAG 正是这样做的\upcite{zhang2025faithfulrag}）。
它采用的是一个\textbf{间接探针}：如果把某条路径加入提示后模型变得更不确定，就推断这条路径与模型的信念相矛盾。
因此，我们把这个方法概括为三步：“\textbf{结构化重写} + \textbf{基于不确定性的路径选择} + \textbf{提示增强}”。
它能否奏效，取决于以下四个隐含假设：
\begin{description}
  \item[A1 外部知识可信。] 冲突时外部知识总是对的。若检索内容本身有误（如被投毒或已过时），方法会主动放大错误。
  \item[A2 熵升高意味着“纠正性”。] 加入路径后熵升高 $\Leftrightarrow$ 路径与参数知识冲突 $\Leftrightarrow$ 路径对答题有用。但无关或嘈杂的信息同样会让模型变得不确定；反之，一条与模型记忆\emph{一致}的正确路径会降低熵，却同样是答题所需的。
  \item[A3 三元组抽取是忠实的。] 抽取由同一个 LLM 完成。面对反事实文本，模型可能用自己的先验“悄悄改正”原文，或干脆漏抽，这样冲突在抽取阶段就已经被抹掉了。
  \item[A4 答案位于查询实体的两跳之内。] 检索只做两跳遍历，而 MuSiQue 等多跳问题可能需要三到四跳。
\end{description}

% ---------------------------------------------------------------------
\subsubsection{三个模块的逐一分析}
\label{sec:cl-modules}

\paragraph{（1）图构建。}
论文的描述是：先对检索内容做语义分段，再用 LLM 逐段抽取三元组 $(h,r,t)$，汇总成图 $\mathcal{G}=(\mathcal{E},\mathcal{R},\mathcal{T})$。
代码中的实现（\texttt{modules.py} 中的 \texttt{generate\_knowledge\_graph}）是：
按 1024 个 token、重叠 128 的窗口切分文本；
对每个分块调用一次抽取提示，再追加一轮“继续抽取”（gleaning），即每个分块固定 2 次 LLM 调用；
输出的实体带有名称、类型和一段自由文本描述，关系带有描述、关键词和权重；
同名实体合并，描述拼接。
这一流程与 LightRAG\upcite{guo2024lightrag} 的图构建实现高度一致。
我们注意到三点：
\begin{itemize}
  \item \textbf{图是“每题局部”的}：每道题都只用这道题自己的检索文本从零建一张图，答完即清空（\texttt{pipeline.py} 中的 \texttt{make\_knowledge\_graph}）。因此它并非一个持久的知识库，建图成本无法在多道题之间摊销，也无法利用其它文档做交叉验证。
  \item \textbf{实体对齐只靠名称完全匹配}（统一大写），别名、指代不会合并，同一实体可能分裂成多个节点。
  \item \textbf{“结构化”是有限度的}：节点和边上挂的都是自然语言描述。最终交给模型的路径文本（见下）大部分仍是自然语言句子，只是按实体和关系组织了起来。
\end{itemize}

\paragraph{（2）图检索。}
论文用式（\ref{eq:cl-ref}）对候选路径打分，并取前 $K$ 条：
\begin{equation}
  \mathrm{Ref}(p)=\alpha\cdot\frac{|\{e\in p\}\cap\mathcal{E}_{imp}|}{|\mathcal{E}_{imp}|}
               +\beta\cdot\frac{|\{r\in p\}\cap\mathcal{R}_{imp}|}{|\mathcal{R}_{imp}|}
  \label{eq:cl-ref}
\end{equation}
代码（\texttt{modules.py} 中的 \texttt{\_build\_query\_context}）的实际流程是：
\begin{enumerate}
  \item 用 LLM 从问题中抽取至多 5 个实体，以及“答案类型”关键词（从图中已有的实体类型里挑选）；
  \item 对每个问题实体，在实体名称向量库中检索最相似的 10 个节点（余弦相似度需大于 0.2），作为起点；
  \item 从每个起点枚举所有两跳路径；找不到路径的孤立起点只保留相似度最高的 20\%；
  \item 统计每条路径上属于“答案类型”的节点个数；
  \item 用整个问题在关系向量库中检索若干条边，其中与关键实体相连的边记为“有效边”；一条路径只要包含至少一条有效边，就被保留；
  \item 删除被其它路径包含的子路径，再把每条路径渲染为“Path / Nodes / Edges”三段文本。
\end{enumerate}
这套逻辑与 MiniRAG\upcite{fan2025minirag} 的异构图检索在函数结构和提示模板上高度一致。
关键在于：第 4、5 步算出的计数\textbf{只用于判断去留，并不用于排序}。最终路径列表按路径长度排列，随后由 \texttt{pipeline.py} 中的 \texttt{run} 直接截取前 30 条送入下一模块。
换言之，式（\ref{eq:cl-ref}）中的 $\alpha$、$\beta$ 和按得分取前 $K$ 条的操作在代码中并不存在，“选哪几条路径”实际上完全由下一步的熵过滤决定。

\paragraph{（3）基于熵的冲突消解。}
论文用式（\ref{eq:cl-entropy}）衡量模型输出的不确定性：
\begin{equation}
  H\big(P(\cdot\mid \text{ctx})\big) = -\frac{1}{L}\sum_{t=1}^{L}\sum_{i=1}^{k} p_i^{(t)}\log p_i^{(t)},
  \qquad
  \Delta H_p = H\big(P(\cdot\mid q,p)\big) - H\big(P(\cdot\mid q)\big),
  \label{eq:cl-entropy}
\end{equation}
并保留 $\Delta H_p>\tau$ 的路径。代码实现（\texttt{modules.py} 中的 \texttt{entropy\_filter}，\texttt{llm/backend.py} 中的 \texttt{cal\_entroy}）有以下细节：
\begin{itemize}
  \item 计算 $H(\cdot\mid q)$ 和 $H(\cdot\mid q,p)$ 时，提示中\textbf{不含原始上下文}，只有问题（和选项）以及该条路径；用贪心解码生成，取每个位置概率最高的 10 个词元，重新归一化后按\textbf{自然对数}计算熵（论文写的是 $\log_2$）。
  \item 生成长度 $L$ 被设为\textbf{金标准答案的 token 数}。也就是说，推理阶段用到了标注信息，存在轻微的标签泄漏。熵是在模型自己生成的前 $L$ 个词元上计算的，而不是在正确答案上。
  \item 每条候选路径都要单独调用一次模型，每道题共需 $1+|P|$ 次调用（$|P|\le 30$）。
\end{itemize}
最关键的差异在于选择规则。代码实际执行的是：
\begin{equation}
  \text{keep}(p) = \big[\Delta H_p \ge 0\big] \;\lor\; \Big[\,H_{aug}(p)>0 \;\land\; \frac{H_{base}}{H_{aug}(p)} \le 10^{\tau}\Big],
  \label{eq:cl-coderule}
\end{equation}
然后按 $\Delta H_p$ 从大到小取前 10 条；若一条都不满足，则保留熵最大的那条。
按式（\ref{eq:cl-coderule}），当 $\tau=1$ 时，一条使熵\emph{下降}的路径只有在熵降到原来的十分之一以下时才会被丢弃；当 $\tau=3$ 时，则要降到千分之一以下。
因此这条规则几乎等价于“\textbf{全部保留，按 $\Delta H$ 排序取前 10}”，$\tau$ 的作用非常弱。
这也为论文附录中“$\tau=1$ 与 $\tau=3$ 结果几乎相同”的现象提供了另一种解释：它未必说明方法对超参数鲁棒，更可能只是 $\tau$ 本身几乎不起作用。
反过来，如果真的按论文写的 $\Delta H_p>\tau$ 执行：在 top-10 重归一化下，每个位置的熵上限为 $\ln 10\approx 2.30$，$\tau=1$ 已是很严格的要求，很多题可能一条路径都留不下来（第~\ref{sec:cl-e1}~小节将验证这一点）。

\paragraph{（4）最终生成。}
论文写作 $\text{Response}=\mathcal{M}(q\oplus\mathcal{P}_{corrective})$，读者容易理解为用筛选后的路径\emph{替换}原始上下文。
代码（\texttt{modules.py} 中的 \texttt{predict\_answer}）实际构造的提示同时包含筛选后的路径（标注为 Facts）\textbf{和完整的原始上下文}（标注为 Context），并要求模型“优先依据 Facts，信息不足时再参考 Context”，以思维链加 JSON 的格式输出。
所以更准确的写法是 $\mathcal{M}(q\oplus\mathcal{P}_{kept}\oplus C)$：方法是在原始上下文之上\textbf{追加}一组被强调的事实，而不是替换上下文。
这一点直接影响对论文消融实验和“上下文长度”统计的解读。

% ---------------------------------------------------------------------
\subsubsection{论文描述与开源实现的对照}
\label{sec:cl-gap}

表~\ref{tab:cl-gap}~汇总了上文发现的差异。需要说明的是，开源代码未必就是论文实验所用的版本，因此我们不把这些差异直接当作论文的错误，而是将它们作为后续实验需要控制或检验的因素。

\begin{table}[htbp]
  \centering
  \small
  \caption{论文描述与开源实现的主要差异}
  \label{tab:cl-gap}
  \begin{tabular}{p{1.6cm}p{3.6cm}p{4.4cm}p{4.2cm}}
    \toprule
    环节 & 论文描述 & 代码实现 & 影响 \\
    \midrule
    路径打分 & 式（\ref{eq:cl-ref}），按 $\mathrm{Ref}(p)$ 取前 $K$ 条 & 无 $\alpha,\beta$；计数只决定去留，按路径长度排序后截取前 30 条 & 检索阶段不做相关性排序，选择完全依赖熵过滤 \\
    熵的计算 & $\log_2$，答案的 token 分布 & top-10 重归一化，自然对数；生成长度 = 金标准答案长度 & 推理时使用了标注信息；阈值的量纲与论文不同 \\
    过滤规则 & 保留 $\Delta H_p>\tau$ & 式（\ref{eq:cl-coderule}），再取前 10 条；为空时保留熵最大者 & 几乎全部保留，$\tau$ 作用很弱 \\
    最终输入 & $q\oplus\mathcal{P}_{corrective}$ & $q\oplus\mathcal{P}_{kept}\oplus C$，提示要求优先依据 Facts & 是“增强”而非“替换” \\
    知识图谱 & 知识图谱 & 每题从零构建的局部图 & 建图成本无法摊销，没有跨文档校验 \\
    评测指标 & 准确率 ACC & 归一化后的\emph{预测}是\emph{金标准}的子串即判对；空预测也判对 & 指标偏宽松；我们同时报告严格的精确匹配（EM） \\
    \bottomrule
  \end{tabular}
\end{table}

% ---------------------------------------------------------------------
\subsubsection{方法层面的优缺点小结}
\label{sec:cl-method-summary}

\paragraph{优点。}
（1）\textbf{定位粒度合适}：三元组把冲突限定在一条边上，路径文本比原文更短、噪声更少，这与论文图 3 中模型对正确答案更自信的观察相符。
（2）\textbf{与模型无关、即插即用}：只要能拿到输出词元的对数概率，就能接入任意 RAG 系统，不需要训练。
（3）\textbf{中间结果可审计}：每条候选路径和它的 $\Delta H$ 都可以记录下来，便于定位错误发生在哪一步。我们正是利用这一点设计了逐级漏斗分析（第~\ref{sec:cl-e3}~小节）。

\paragraph{缺点。}
（1）\textbf{成本高}：每道题需要“$2\times$ 分块数（建图）$+1$（关键词）$+1+|P|$（熵）$+1$（生成）”次调用。论文附录给出的单题耗时为 36--62 秒，而标准 RAG 仅需 0.4--2.6 秒。
（2）\textbf{核心信号缺乏直接验证}：A2 是整个方法的立足点，论文只给出了最终准确率，没有直接检验“熵升高的路径是否真的包含正确答案”。
（3）\textbf{对错误外部知识没有防御}：在 A1 不成立时，方法会把错误信息以“Facts”的形式加以强调。
（4）\textbf{误差逐级传递}：抽取漏掉或“改正”的事实，后续模块无从恢复（A3、A4）。
（5）\textbf{依赖标注与对数概率}：生成长度用到了金标准答案；而很多闭源 API 并不提供对数概率。

\paragraph{待检验的假设。}
基于以上分析，我们提出四个可检验的假设，并与后续实验一一对应，如表~\ref{tab:cl-hypo}~所示。

\begin{table}[htbp]
  \centering
  \small
  \caption{由方法分析得到的假设与对应实验}
  \label{tab:cl-hypo}
  \begin{tabular}{p{0.6cm}p{8.6cm}p{2.0cm}p{1.6cm}}
    \toprule
    编号 & 假设 & 对应前提 & 实验 \\
    \midrule
    H1 & $\Delta H$ 区分“含答案路径”与“不含答案路径”的能力有限 & A2 & E1 \\
    H2 & 按论文规则过滤会丢掉大量有用路径；代码规则的收益主要来自排序与提示增强，而非“过滤冲突” & A2、实现差异 & E2 \\
    H3 & 答案主要在抽取、检索阶段丢失，而非在过滤阶段 & A3、A4 & E3 \\
    H4 & 当外部知识本身有误时，方法比标准 RAG 更容易盲从 & A1 & R3 \\
    \bottomrule
  \end{tabular}
\end{table}
                 % X.2   Sean
% 复现实验 —— 负责人：Hakimi
% （占位文件：由 Hakimi 撰写。标签 sec:cl-repro 已被其它小节引用，请保留。）
\subsection{复现实验}
\label{sec:cl-repro}

\subsubsection{实验设置}
\cltodo{Hakimi：子集抽样方式（seed=42）、模型、运行环境、对代码所做的改动（逐题记录、并发）、成本}

\subsubsection{结果与对比}
\cltodo{Hakimi：R1、R2 结果，并与论文表 1、表 2 对比}

\subsubsection{原文实验设计评述}
\cltodo{Hakimi：显著性检验方式、CPR 定义、Qwen-72B 上的异常值、RealtimeQA 样本量}
           % X.3   Hakimi

\subsection{探究实验}
\label{sec:cl-probe}
本节的实验围绕第~\ref{sec:cl-method-summary}~小节提出的四个假设展开：
\ref{sec:cl-e1}、\ref{sec:cl-e3}~小节检验方法自身的设计假设，
\ref{sec:cl-r3}、\ref{sec:cl-r4}~小节检验论文实验结论的可靠性与可推广性。
% X.4.1–X.4.2 探究实验（方法假设）—— 负责人：Sean
% 数据来源：results/trace_*.jsonl（Hakimi 运行 R1/R2 时产生）
% 脚本：experiments/analysis/e1_entropy.py, e2_filters.py, e3_funnel.py
% 图：figure/cl_method/ ；表：experiments/out/cl_method/

\paragraph{共同设置。}
本节实验全部复用复现实验（第~\ref{sec:cl-repro}~小节）中 TruthfulRAG 的运行记录。
我们在开源代码中加入了逐题记录，保存每道题的知识图谱规模、全部候选路径及其熵值 $H$ 与 $\Delta H$、被保留的路径编号以及最终答案。
记录逻辑只在原有计算之外附加输出，不改变方法的行为。我们还按式（\ref{eq:cl-coderule}）独立重写了一遍选择规则，在离线测试中，它选出的路径与原实现记录的路径完全一致；E2 中的“代码规则”正是基于这一重写版本。
一条路径“含答案”的判定方法是：按官方评测脚本的规则归一化后，金标准答案作为完整词组出现在路径文本中。
对反事实数据集（MuSiQue、SQuAD），金标准是检索文本中被篡改后的答案，也就是方法应当“跟随”的那个答案。
模型为 GPT-4o-mini，每个数据集抽取样本 \cltodo{N}~条。

% ---------------------------------------------------------------------
\subsubsection{熵过滤机制是否成立（E1、E2）}
\label{sec:cl-e1}

\paragraph{E1：$\Delta H$ 能否识别“纠正性路径”。}
如果假设 A2 成立，含答案的路径应当系统性地具有更大的 $\Delta H$。
我们把所有候选路径按“是否含答案”分为两组，比较两组 $\Delta H$ 的分布，并报告三个指标：
（1）以 $\Delta H$ 为打分、区分两组的 AUC（0.5 即等同于随机）；
（2）单侧 Mann--Whitney U 检验的 $p$ 值；
（3）在至少有一条含答案路径的题目中，$\Delta H$ 最大的那条路径恰好含答案的比例（Hit@1），并与随机排序下的期望值对比。

\begin{figure}[htbp]
  \centering
  \clfig{figure/cl_method/e1_dh_by_label.pdf}{0.85\textwidth}
  \caption{含答案路径与不含答案路径的 $\Delta H$ 分布（横线为中位数，虚线为 $\Delta H=0$）}
  \label{fig:cl-e1}
\end{figure}

\cltodo{表：各数据集的路径数、含答案比例、两组平均 $\Delta H$、AUC、$p$ 值、Hit@1（来自 \texttt{e1\_entropy.tex}）}

\cltodo{结果分析：AUC 是否显著高于 0.5；含答案路径中 $\Delta H>0$ 的比例；在 golden（无冲突）数据上，含答案路径的 $\Delta H$ 是否转为负值——若是，则说明熵升高反映的是“与记忆冲突”，而不是“对答题有用”，从而支持 H1}

\paragraph{E2：过滤策略对比。}
候选路径及其熵值均已记录，因此更换过滤规则时只需重做最后一步生成（每题 1 次调用）。我们比较以下策略：
\begin{itemize}
  \item \textbf{代码规则}：式（\ref{eq:cl-coderule}），即原始实现；
  \item \textbf{论文规则}：只保留 $\Delta H_p>\tau$ 的路径；
  \item \textbf{全部保留}：不过滤，对应代码自带的消融设置；
  \item \textbf{随机}：随机保留与代码规则同样数量的路径；
  \item \textbf{反向}：保留 $\Delta H$ 最\emph{小}的同样数量的路径；
  \item \textbf{不加路径}：只用问题和原始上下文，提示格式不变。
\end{itemize}
其中，“随机”和“反向”用来检验 $\Delta H$ 排序本身是否有价值；“不加路径”用来衡量整个知识图谱分支的净贡献。
我们还在 $\tau\in\{0,0.5,1,2,3\}$ 上扫描论文规则，统计“保留路径中含答案”和“一条路径都没保留”的题目比例。

\begin{figure}[htbp]
  \centering
  \begin{minipage}{0.48\textwidth}\centering
    \clfig{figure/cl_method/e2_accuracy.pdf}{\linewidth}
  \end{minipage}\hfill
  \begin{minipage}{0.48\textwidth}\centering
    \clfig{figure/cl_method/e2_tau_coverage.pdf}{\linewidth}
  \end{minipage}
  \caption{左：不同过滤策略的平均准确率；右：论文规则下阈值 $\tau$ 对保留结果的影响}
  \label{fig:cl-e2}
\end{figure}

\cltodo{表：各策略在各数据集上的准确率（来自 \texttt{e2\_accuracy.tex}）}

\cltodo{结果分析：代码规则与“随机”、“反向”、“不加路径”之间的差距；论文规则在 $\tau=1$ 时空集的比例；据此判断 H2}

% ---------------------------------------------------------------------
\subsubsection{流水线逐级损失与错误归因（E3）}
\label{sec:cl-e3}

为定位答案在哪一步丢失，我们逐题检查答案是否依次出现在：
S1 原始检索文本、S2 构建出的知识图谱、S3 候选路径、S4 被保留的路径中，以及 S5 最终是否答对。
由于最终提示中仍包含完整的原始上下文（第~\ref{sec:cl-modules}~小节），S5 并不受 S4 约束。
因此我们额外比较“保留路径含答案”与“保留路径不含答案”两种情况下的准确率，以区分模型的正确答案究竟来自路径，还是来自原文。

\begin{figure}[htbp]
  \centering
  \clfig{figure/cl_method/e3_funnel.pdf}{\textwidth}
  \caption{各数据集上答案的逐级保留率（虚线为标准 RAG 的准确率）}
  \label{fig:cl-e3}
\end{figure}

\cltodo{结果分析：最大的一级损失发生在哪里（S1→S2 抽取，还是 S2→S3 检索）；反事实数据上 S1→S2 的损失是否明显高于 golden 数据（即 A3 中的“抽取时被模型先验改正”）；S4 不成立时的准确率}

\cltodo{（P1）人工归因：从答错的样本中抽取 30 条，归入抽取错误、检索遗漏、过滤误删、生成错误四类，并给出一个典型案例}
          % X.4.1–X.4.2  Sean
% 补充实验 R3–R4—— 负责人：Hakimi
% （占位文件：由 Hakimi 撰写。标签 sec:cl-r3 / sec:cl-r4 已被其它小节引用，请保留。）
\subsubsection{外部知识错误下的鲁棒性}
\label{sec:cl-r3}
\cltodo{Hakimi：用 negative 数据的记录、以 gold\_true 重新评估；检验假设 H4（表~\ref{tab:cl-hypo}）}

\subsubsection{统计显著性与计算成本}
\label{sec:cl-r4}
\cltodo{Hakimi：McNemar 检验、bootstrap 置信区间、准确率与成本的关系}
      % X.4.3–X.4.4  Hakimi

\subsection{综合讨论：局限与改进方向}
\label{sec:cl-discussion}
% X.5 前半：方法层面的局限与改进 —— 负责人：Sean（D2 撰写）
\cltodo{Sean：结合 E1–E3 的结果，讨论方法层面的局限，并提出改进方向（例如：用“答案分布是否翻转”替代熵变化作为冲突信号；先在原文中定位冲突句、只抽取相关三元组以降低建图成本；对外部知识增加可信度校验以防御 A1 失效）}
     % X.5 前半  Sean
% 局限与展望：实验层面 —— 负责人：Hakimi
\subsubsection{实验层面}
\label{sec:cl-discussion-exp}
\cltodo{Hakimi：实验层面的局限（样本规模、单一模型、评测指标等）}
 % X.5 后半  Hakimi

% 小结 —— 负责人：Hakimi
\subsection{小结}
\label{sec:cl-summary}
\cltodo{Hakimi：半页小结}
                % X.6   Hakimi

% 每个小节一个文件，文件头注明负责人；只改自己负责的文件。
% =====================================================================
\providecommand{\cltodo}[1]{\textcolor{red}{[待补：#1]}}
% 图片若尚未生成，则显示占位框，保证随时可编译
\providecommand{\clfig}[2]{\IfFileExists{#1}{\includegraphics[width=#2]{#1}}{\fbox{\parbox{#2}{\centering\cltodo{图片 \texttt{\detokenize{#1}} 尚未生成}}}}}

\section{TruthfulRAG：基于知识图谱的 RAG 事实级冲突消解}
\label{sec:cl}

% 论文概述 —— 负责人：Sean
\subsection{论文概述}
\label{sec:cl-overview}

本节依次概述论文的研究背景、问题定义、方法框架以及实验设置与结论，内容均依据论文原文整理。
本组对方法的分析与评价集中在第~\ref{sec:cl-method}~节。

% ---------------------------------------------------------------------
\subsubsection{研究背景}

检索增强生成（Retrieval-Augmented Generation, RAG）将外部检索得到的文档拼接进提示，用以弥补大语言模型参数知识过时或缺失的不足\upcite{lewis2020retrieval,gao2023retrieval}。
随着外部知识库不断更新，检索内容与模型在预训练阶段记住的知识可能互相矛盾，这一现象称为知识冲突。
已有研究表明，发生冲突时模型往往倾向于坚持自身的记忆而忽视检索证据\upcite{longpre2021entity,xie2023adaptive}，从而降低 RAG 系统输出的准确性。

论文将已有的冲突消解方法归为两类。
词元级方法在解码阶段调整输出词元的概率分布，例如对比有无上下文时的输出分布\upcite{yuan2024discerning,jin2024tug}；
语义级方法在文段或陈述层面对齐内部与外部知识，例如 FaithfulRAG 先让模型显式写出自身的参数知识，再与检索内容比对\upcite{zhang2025faithfulrag}。
作者认为两类方法依赖的表示较为零散，难以准确刻画事实之间的关联，因而无法精确定位发生冲突的具体事实。

% ---------------------------------------------------------------------
\subsubsection{问题定义}

记 $\mathcal{M}$ 为具有参数知识 $\mathcal{K}_p$ 的大语言模型。
给定问题 $q$，RAG 系统从外部知识库 $\mathcal{K}_e$ 中检索文档，并以问题和检索内容为条件生成回答。
若存在与 $q$ 相关的两条事实陈述 $f_p\in\mathcal{K}_p$ 与 $f_e\in\mathcal{K}_e$，且二者在事实上不一致，则称发生了知识冲突。
论文采取外部知识优先的立场，目标是在冲突发生时使模型遵循检索内容作答，同时在无冲突时不损害原有性能。
围绕这一目标，论文提出了三个技术问题：如何表示事实知识以便检测冲突，如何为给定问题检索并排序相关事实，以及如何使模型在面对冲突证据时做出可靠决策。

% ---------------------------------------------------------------------
\subsubsection{方法框架}

TruthfulRAG 由图构建、图检索和冲突消解三个模块串联而成，整体流程如图~\ref{fig:cl-pipeline}~所示。

\begin{figure}[tbp]
  \centering
  \clfig{figure/cl_method/pipeline.png}{0.92\textwidth}
  \caption{TruthfulRAG 的整体流程，依次为图构建、图检索与基于熵的冲突消解（取自论文开源仓库）}
  \label{fig:cl-pipeline}
\end{figure}

图构建模块将检索内容划分为语义连贯的文段，由模型 $\mathcal{M}$ 从每个文段中抽取三元组 $(h,r,t)$，汇总得到知识图谱 $\mathcal{G}=(\mathcal{E},\mathcal{R},\mathcal{T})$，其中 $\mathcal{E}$、$\mathcal{R}$、$\mathcal{T}$ 分别为实体、关系与三元组集合。

图检索模块先从问题中抽取关键要素，按语义相似度选出最相关的实体集合 $\mathcal{E}_{imp}$ 与关系集合 $\mathcal{R}_{imp}$，再从每个关键实体出发做两跳遍历，得到初始路径集合 $\mathcal{P}_{init}$。
每条路径按其对关键实体与关系的覆盖程度打分：
\begin{equation}
  \mathrm{Ref}(p)=\alpha\,\frac{|\{e\in p\}\cap\mathcal{E}_{imp}|}{|\mathcal{E}_{imp}|}
               +\beta\,\frac{|\{r\in p\}\cap\mathcal{R}_{imp}|}{|\mathcal{R}_{imp}|},
  \label{eq:cl-ref}
\end{equation}
得分最高的 $K$ 条路径构成 $\mathcal{P}_{super}$。每条路径被展开为路径序列、实体属性与关系属性三部分文本。

冲突消解模块以输出熵度量模型的不确定性。设生成答案长度为 $L$，第 $t$ 个位置上前 $k$ 个候选词元的概率为 $p_i^{(t)}$，则
\begin{equation}
  H = -\frac{1}{L}\sum_{t=1}^{L}\sum_{i=1}^{k} p_i^{(t)}\log_2 p_i^{(t)} .
  \label{eq:cl-entropy}
\end{equation}
分别计算只给出问题时的熵 $H_{param}$ 与给出问题和路径 $p$ 时的熵 $H_{aug}(p)$，二者之差 $\Delta H_p=H_{aug}(p)-H_{param}$ 被视为冲突的指示量。
论文认为 $\Delta H_p$ 为正说明该路径与模型的参数知识相矛盾，对模型具有纠正作用，因此保留
\begin{equation}
  \mathcal{P}_{corrective}=\{\,p\in\mathcal{P}_{super}\mid \Delta H_p>\tau\,\},
  \label{eq:cl-paperrule}
\end{equation}
并以 $\mathcal{M}(q\oplus\mathcal{P}_{corrective})$ 生成最终回答，其中 $\tau$ 为阈值，$\oplus$ 表示文本拼接。

% ---------------------------------------------------------------------
\subsubsection{实验设置与主要结论}

论文在四个数据集上评测：FaithEval\upcite{ming2025faitheval} 包含反事实与不一致的上下文；MuSiQue\upcite{trivedi2022musique} 与 SQuAD\upcite{rajpurkar2016squad} 采用 KRE\upcite{ying2024intuitive} 构造的反事实版本；RealtimeQA\upcite{kasai2023realtime} 侧重时效性冲突。
所用模型为 GPT-4o-mini、Qwen2.5-7B-Instruct 与 Mistral-7B-Instruct。
对比基线包括不使用检索的直接生成、标准 RAG、提示优化方法 KRE、解码干预方法 COIECD\upcite{yuan2024discerning} 以及 FaithfulRAG，评价指标为准确率。

论文报告的主要结论有四点。
其一，TruthfulRAG 在三个模型上均取得最高的平均准确率，在 GPT-4o-mini、Qwen2.5-7B 与 Mistral-7B 上分别为 78.8\%、78.3\% 与 81.3\%，次优的 FaithfulRAG 分别为 76.5\%、78.1\% 与 80.7\%。
其二，在 MuSiQue 与 SQuAD 的无冲突版本上，该方法的准确率不低于标准 RAG。
其三，与自然语言上下文相比，结构化的推理路径使模型赋予正确答案更高的对数概率。
其四，消融实验表明图模块与冲突消解模块均对性能有贡献。
               % X.1   Sean
% X.2 方法剖析 —— 负责人：Sean
\subsection{方法剖析}
\label{sec:cl-method}

本节的分析同时依据论文正文和作者开源的代码。我们逐函数阅读了实现，下文提到的代码位置均以“文件名中的函数名”标注，便于核对。
分析按“问题形式化 → 三个模块 → 论文与实现的对照 → 优缺点与待检验假设”的顺序展开。

% ---------------------------------------------------------------------
\subsubsection{问题形式化与设计动机}
\label{sec:cl-formal}

记 LLM 为 $\mathcal{M}$，问题为 $q$，检索到的上下文为 $C$。
模型只凭参数知识作答时的输出分布为 $p_{\mathcal{M}}(y\mid q)$，给定上下文后为 $p_{\mathcal{M}}(y\mid q, C)$。
令 $a_{\text{param}}=\arg\max_y p_{\mathcal{M}}(y\mid q)$ 为模型“记住的”答案，$a_{\text{ctx}}$ 为上下文 $C$ 支持的答案。
\textbf{知识冲突}即 $a_{\text{param}}\neq a_{\text{ctx}}$ 的情形。论文采取“外部知识优先”的立场，目标是让最终输出等于 $a_{\text{ctx}}$。
这一立场隐含了一个前提：外部知识是对的。我们在下文把它记作假设~A1。

\paragraph{为什么是“事实级”。}
我们对这一动机的理解是：冲突往往只集中在\textbf{极少数}事实上。
以论文附录的例子为例，检索文本把“新拉雷多市所在的州”从 Tamaulipas 改成了 Sinaloa，而同一段文字中其余几十条事实都与模型的记忆一致。
在整段文本上做“语义级”比较时，这一条冲突会被大量一致的信息稀释；在输出词元上做“词元级”调整时，方法又并不知道是哪条事实引起了分歧。
若把文本拆成三元组，冲突就被定位到一条边（Nuevo Laredo, located\_in, Sinaloa）上，可以单独处理。

\paragraph{方法的实质。}
值得注意的是，TruthfulRAG 并\textbf{没有}显式地把模型的参数知识写成三元组、再与外部三元组逐条比对（FaithfulRAG 正是这样做的\upcite{zhang2025faithfulrag}）。
它采用的是一个\textbf{间接探针}：如果把某条路径加入提示后模型变得更不确定，就推断这条路径与模型的信念相矛盾。
因此，我们把这个方法概括为三步：“\textbf{结构化重写} + \textbf{基于不确定性的路径选择} + \textbf{提示增强}”。
它能否奏效，取决于以下四个隐含假设：
\begin{description}
  \item[A1 外部知识可信。] 冲突时外部知识总是对的。若检索内容本身有误（如被投毒或已过时），方法会主动放大错误。
  \item[A2 熵升高意味着“纠正性”。] 加入路径后熵升高 $\Leftrightarrow$ 路径与参数知识冲突 $\Leftrightarrow$ 路径对答题有用。但无关或嘈杂的信息同样会让模型变得不确定；反之，一条与模型记忆\emph{一致}的正确路径会降低熵，却同样是答题所需的。
  \item[A3 三元组抽取是忠实的。] 抽取由同一个 LLM 完成。面对反事实文本，模型可能用自己的先验“悄悄改正”原文，或干脆漏抽，这样冲突在抽取阶段就已经被抹掉了。
  \item[A4 答案位于查询实体的两跳之内。] 检索只做两跳遍历，而 MuSiQue 等多跳问题可能需要三到四跳。
\end{description}

% ---------------------------------------------------------------------
\subsubsection{三个模块的逐一分析}
\label{sec:cl-modules}

\paragraph{（1）图构建。}
论文的描述是：先对检索内容做语义分段，再用 LLM 逐段抽取三元组 $(h,r,t)$，汇总成图 $\mathcal{G}=(\mathcal{E},\mathcal{R},\mathcal{T})$。
代码中的实现（\texttt{modules.py} 中的 \texttt{generate\_knowledge\_graph}）是：
按 1024 个 token、重叠 128 的窗口切分文本；
对每个分块调用一次抽取提示，再追加一轮“继续抽取”（gleaning），即每个分块固定 2 次 LLM 调用；
输出的实体带有名称、类型和一段自由文本描述，关系带有描述、关键词和权重；
同名实体合并，描述拼接。
这一流程与 LightRAG\upcite{guo2024lightrag} 的图构建实现高度一致。
我们注意到三点：
\begin{itemize}
  \item \textbf{图是“每题局部”的}：每道题都只用这道题自己的检索文本从零建一张图，答完即清空（\texttt{pipeline.py} 中的 \texttt{make\_knowledge\_graph}）。因此它并非一个持久的知识库，建图成本无法在多道题之间摊销，也无法利用其它文档做交叉验证。
  \item \textbf{实体对齐只靠名称完全匹配}（统一大写），别名、指代不会合并，同一实体可能分裂成多个节点。
  \item \textbf{“结构化”是有限度的}：节点和边上挂的都是自然语言描述。最终交给模型的路径文本（见下）大部分仍是自然语言句子，只是按实体和关系组织了起来。
\end{itemize}

\paragraph{（2）图检索。}
论文用式（\ref{eq:cl-ref}）对候选路径打分，并取前 $K$ 条：
\begin{equation}
  \mathrm{Ref}(p)=\alpha\cdot\frac{|\{e\in p\}\cap\mathcal{E}_{imp}|}{|\mathcal{E}_{imp}|}
               +\beta\cdot\frac{|\{r\in p\}\cap\mathcal{R}_{imp}|}{|\mathcal{R}_{imp}|}
  \label{eq:cl-ref}
\end{equation}
代码（\texttt{modules.py} 中的 \texttt{\_build\_query\_context}）的实际流程是：
\begin{enumerate}
  \item 用 LLM 从问题中抽取至多 5 个实体，以及“答案类型”关键词（从图中已有的实体类型里挑选）；
  \item 对每个问题实体，在实体名称向量库中检索最相似的 10 个节点（余弦相似度需大于 0.2），作为起点；
  \item 从每个起点枚举所有两跳路径；找不到路径的孤立起点只保留相似度最高的 20\%；
  \item 统计每条路径上属于“答案类型”的节点个数；
  \item 用整个问题在关系向量库中检索若干条边，其中与关键实体相连的边记为“有效边”；一条路径只要包含至少一条有效边，就被保留；
  \item 删除被其它路径包含的子路径，再把每条路径渲染为“Path / Nodes / Edges”三段文本。
\end{enumerate}
这套逻辑与 MiniRAG\upcite{fan2025minirag} 的异构图检索在函数结构和提示模板上高度一致。
关键在于：第 4、5 步算出的计数\textbf{只用于判断去留，并不用于排序}。最终路径列表按路径长度排列，随后由 \texttt{pipeline.py} 中的 \texttt{run} 直接截取前 30 条送入下一模块。
换言之，式（\ref{eq:cl-ref}）中的 $\alpha$、$\beta$ 和按得分取前 $K$ 条的操作在代码中并不存在，“选哪几条路径”实际上完全由下一步的熵过滤决定。

\paragraph{（3）基于熵的冲突消解。}
论文用式（\ref{eq:cl-entropy}）衡量模型输出的不确定性：
\begin{equation}
  H\big(P(\cdot\mid \text{ctx})\big) = -\frac{1}{L}\sum_{t=1}^{L}\sum_{i=1}^{k} p_i^{(t)}\log p_i^{(t)},
  \qquad
  \Delta H_p = H\big(P(\cdot\mid q,p)\big) - H\big(P(\cdot\mid q)\big),
  \label{eq:cl-entropy}
\end{equation}
并保留 $\Delta H_p>\tau$ 的路径。代码实现（\texttt{modules.py} 中的 \texttt{entropy\_filter}，\texttt{llm/backend.py} 中的 \texttt{cal\_entroy}）有以下细节：
\begin{itemize}
  \item 计算 $H(\cdot\mid q)$ 和 $H(\cdot\mid q,p)$ 时，提示中\textbf{不含原始上下文}，只有问题（和选项）以及该条路径；用贪心解码生成，取每个位置概率最高的 10 个词元，重新归一化后按\textbf{自然对数}计算熵（论文写的是 $\log_2$）。
  \item 生成长度 $L$ 被设为\textbf{金标准答案的 token 数}。也就是说，推理阶段用到了标注信息，存在轻微的标签泄漏。熵是在模型自己生成的前 $L$ 个词元上计算的，而不是在正确答案上。
  \item 每条候选路径都要单独调用一次模型，每道题共需 $1+|P|$ 次调用（$|P|\le 30$）。
\end{itemize}
最关键的差异在于选择规则。代码实际执行的是：
\begin{equation}
  \text{keep}(p) = \big[\Delta H_p \ge 0\big] \;\lor\; \Big[\,H_{aug}(p)>0 \;\land\; \frac{H_{base}}{H_{aug}(p)} \le 10^{\tau}\Big],
  \label{eq:cl-coderule}
\end{equation}
然后按 $\Delta H_p$ 从大到小取前 10 条；若一条都不满足，则保留熵最大的那条。
按式（\ref{eq:cl-coderule}），当 $\tau=1$ 时，一条使熵\emph{下降}的路径只有在熵降到原来的十分之一以下时才会被丢弃；当 $\tau=3$ 时，则要降到千分之一以下。
因此这条规则几乎等价于“\textbf{全部保留，按 $\Delta H$ 排序取前 10}”，$\tau$ 的作用非常弱。
这也为论文附录中“$\tau=1$ 与 $\tau=3$ 结果几乎相同”的现象提供了另一种解释：它未必说明方法对超参数鲁棒，更可能只是 $\tau$ 本身几乎不起作用。
反过来，如果真的按论文写的 $\Delta H_p>\tau$ 执行：在 top-10 重归一化下，每个位置的熵上限为 $\ln 10\approx 2.30$，$\tau=1$ 已是很严格的要求，很多题可能一条路径都留不下来（第~\ref{sec:cl-e1}~小节将验证这一点）。

\paragraph{（4）最终生成。}
论文写作 $\text{Response}=\mathcal{M}(q\oplus\mathcal{P}_{corrective})$，读者容易理解为用筛选后的路径\emph{替换}原始上下文。
代码（\texttt{modules.py} 中的 \texttt{predict\_answer}）实际构造的提示同时包含筛选后的路径（标注为 Facts）\textbf{和完整的原始上下文}（标注为 Context），并要求模型“优先依据 Facts，信息不足时再参考 Context”，以思维链加 JSON 的格式输出。
所以更准确的写法是 $\mathcal{M}(q\oplus\mathcal{P}_{kept}\oplus C)$：方法是在原始上下文之上\textbf{追加}一组被强调的事实，而不是替换上下文。
这一点直接影响对论文消融实验和“上下文长度”统计的解读。

% ---------------------------------------------------------------------
\subsubsection{论文描述与开源实现的对照}
\label{sec:cl-gap}

表~\ref{tab:cl-gap}~汇总了上文发现的差异。需要说明的是，开源代码未必就是论文实验所用的版本，因此我们不把这些差异直接当作论文的错误，而是将它们作为后续实验需要控制或检验的因素。

\begin{table}[htbp]
  \centering
  \small
  \caption{论文描述与开源实现的主要差异}
  \label{tab:cl-gap}
  \begin{tabular}{p{1.6cm}p{3.6cm}p{4.4cm}p{4.2cm}}
    \toprule
    环节 & 论文描述 & 代码实现 & 影响 \\
    \midrule
    路径打分 & 式（\ref{eq:cl-ref}），按 $\mathrm{Ref}(p)$ 取前 $K$ 条 & 无 $\alpha,\beta$；计数只决定去留，按路径长度排序后截取前 30 条 & 检索阶段不做相关性排序，选择完全依赖熵过滤 \\
    熵的计算 & $\log_2$，答案的 token 分布 & top-10 重归一化，自然对数；生成长度 = 金标准答案长度 & 推理时使用了标注信息；阈值的量纲与论文不同 \\
    过滤规则 & 保留 $\Delta H_p>\tau$ & 式（\ref{eq:cl-coderule}），再取前 10 条；为空时保留熵最大者 & 几乎全部保留，$\tau$ 作用很弱 \\
    最终输入 & $q\oplus\mathcal{P}_{corrective}$ & $q\oplus\mathcal{P}_{kept}\oplus C$，提示要求优先依据 Facts & 是“增强”而非“替换” \\
    知识图谱 & 知识图谱 & 每题从零构建的局部图 & 建图成本无法摊销，没有跨文档校验 \\
    评测指标 & 准确率 ACC & 归一化后的\emph{预测}是\emph{金标准}的子串即判对；空预测也判对 & 指标偏宽松；我们同时报告严格的精确匹配（EM） \\
    \bottomrule
  \end{tabular}
\end{table}

% ---------------------------------------------------------------------
\subsubsection{方法层面的优缺点小结}
\label{sec:cl-method-summary}

\paragraph{优点。}
（1）\textbf{定位粒度合适}：三元组把冲突限定在一条边上，路径文本比原文更短、噪声更少，这与论文图 3 中模型对正确答案更自信的观察相符。
（2）\textbf{与模型无关、即插即用}：只要能拿到输出词元的对数概率，就能接入任意 RAG 系统，不需要训练。
（3）\textbf{中间结果可审计}：每条候选路径和它的 $\Delta H$ 都可以记录下来，便于定位错误发生在哪一步。我们正是利用这一点设计了逐级漏斗分析（第~\ref{sec:cl-e3}~小节）。

\paragraph{缺点。}
（1）\textbf{成本高}：每道题需要“$2\times$ 分块数（建图）$+1$（关键词）$+1+|P|$（熵）$+1$（生成）”次调用。论文附录给出的单题耗时为 36--62 秒，而标准 RAG 仅需 0.4--2.6 秒。
（2）\textbf{核心信号缺乏直接验证}：A2 是整个方法的立足点，论文只给出了最终准确率，没有直接检验“熵升高的路径是否真的包含正确答案”。
（3）\textbf{对错误外部知识没有防御}：在 A1 不成立时，方法会把错误信息以“Facts”的形式加以强调。
（4）\textbf{误差逐级传递}：抽取漏掉或“改正”的事实，后续模块无从恢复（A3、A4）。
（5）\textbf{依赖标注与对数概率}：生成长度用到了金标准答案；而很多闭源 API 并不提供对数概率。

\paragraph{待检验的假设。}
基于以上分析，我们提出四个可检验的假设，并与后续实验一一对应，如表~\ref{tab:cl-hypo}~所示。

\begin{table}[htbp]
  \centering
  \small
  \caption{由方法分析得到的假设与对应实验}
  \label{tab:cl-hypo}
  \begin{tabular}{p{0.6cm}p{8.6cm}p{2.0cm}p{1.6cm}}
    \toprule
    编号 & 假设 & 对应前提 & 实验 \\
    \midrule
    H1 & $\Delta H$ 区分“含答案路径”与“不含答案路径”的能力有限 & A2 & E1 \\
    H2 & 按论文规则过滤会丢掉大量有用路径；代码规则的收益主要来自排序与提示增强，而非“过滤冲突” & A2、实现差异 & E2 \\
    H3 & 答案主要在抽取、检索阶段丢失，而非在过滤阶段 & A3、A4 & E3 \\
    H4 & 当外部知识本身有误时，方法比标准 RAG 更容易盲从 & A1 & R3 \\
    \bottomrule
  \end{tabular}
\end{table}
                 % X.2   Sean
% 复现实验 —— 负责人：Hakimi
% （占位文件：由 Hakimi 撰写。标签 sec:cl-repro 已被其它小节引用，请保留。）
\subsection{复现实验}
\label{sec:cl-repro}

\subsubsection{实验设置}
\cltodo{Hakimi：子集抽样方式（seed=42）、模型、运行环境、对代码所做的改动（逐题记录、并发）、成本}

\subsubsection{结果与对比}
\cltodo{Hakimi：R1、R2 结果，并与论文表 1、表 2 对比}

\subsubsection{原文实验设计评述}
\cltodo{Hakimi：显著性检验方式、CPR 定义、Qwen-72B 上的异常值、RealtimeQA 样本量}
           % X.3   Hakimi

\subsection{探究实验}
\label{sec:cl-probe}
本节的实验围绕第~\ref{sec:cl-method-summary}~小节提出的四个假设展开：
\ref{sec:cl-e1}、\ref{sec:cl-e3}~小节检验方法自身的设计假设，
\ref{sec:cl-r3}、\ref{sec:cl-r4}~小节检验论文实验结论的可靠性与可推广性。
% X.4.1–X.4.2 探究实验（方法假设）—— 负责人：Sean
% 数据来源：results/trace_*.jsonl（Hakimi 运行 R1/R2 时产生）
% 脚本：experiments/analysis/e1_entropy.py, e2_filters.py, e3_funnel.py
% 图：figure/cl_method/ ；表：experiments/out/cl_method/

\paragraph{共同设置。}
本节实验全部复用复现实验（第~\ref{sec:cl-repro}~小节）中 TruthfulRAG 的运行记录。
我们在开源代码中加入了逐题记录，保存每道题的知识图谱规模、全部候选路径及其熵值 $H$ 与 $\Delta H$、被保留的路径编号以及最终答案。
记录逻辑只在原有计算之外附加输出，不改变方法的行为。我们还按式（\ref{eq:cl-coderule}）独立重写了一遍选择规则，在离线测试中，它选出的路径与原实现记录的路径完全一致；E2 中的“代码规则”正是基于这一重写版本。
一条路径“含答案”的判定方法是：按官方评测脚本的规则归一化后，金标准答案作为完整词组出现在路径文本中。
对反事实数据集（MuSiQue、SQuAD），金标准是检索文本中被篡改后的答案，也就是方法应当“跟随”的那个答案。
模型为 GPT-4o-mini，每个数据集抽取样本 \cltodo{N}~条。

% ---------------------------------------------------------------------
\subsubsection{熵过滤机制是否成立（E1、E2）}
\label{sec:cl-e1}

\paragraph{E1：$\Delta H$ 能否识别“纠正性路径”。}
如果假设 A2 成立，含答案的路径应当系统性地具有更大的 $\Delta H$。
我们把所有候选路径按“是否含答案”分为两组，比较两组 $\Delta H$ 的分布，并报告三个指标：
（1）以 $\Delta H$ 为打分、区分两组的 AUC（0.5 即等同于随机）；
（2）单侧 Mann--Whitney U 检验的 $p$ 值；
（3）在至少有一条含答案路径的题目中，$\Delta H$ 最大的那条路径恰好含答案的比例（Hit@1），并与随机排序下的期望值对比。

\begin{figure}[htbp]
  \centering
  \clfig{figure/cl_method/e1_dh_by_label.pdf}{0.85\textwidth}
  \caption{含答案路径与不含答案路径的 $\Delta H$ 分布（横线为中位数，虚线为 $\Delta H=0$）}
  \label{fig:cl-e1}
\end{figure}

\cltodo{表：各数据集的路径数、含答案比例、两组平均 $\Delta H$、AUC、$p$ 值、Hit@1（来自 \texttt{e1\_entropy.tex}）}

\cltodo{结果分析：AUC 是否显著高于 0.5；含答案路径中 $\Delta H>0$ 的比例；在 golden（无冲突）数据上，含答案路径的 $\Delta H$ 是否转为负值——若是，则说明熵升高反映的是“与记忆冲突”，而不是“对答题有用”，从而支持 H1}

\paragraph{E2：过滤策略对比。}
候选路径及其熵值均已记录，因此更换过滤规则时只需重做最后一步生成（每题 1 次调用）。我们比较以下策略：
\begin{itemize}
  \item \textbf{代码规则}：式（\ref{eq:cl-coderule}），即原始实现；
  \item \textbf{论文规则}：只保留 $\Delta H_p>\tau$ 的路径；
  \item \textbf{全部保留}：不过滤，对应代码自带的消融设置；
  \item \textbf{随机}：随机保留与代码规则同样数量的路径；
  \item \textbf{反向}：保留 $\Delta H$ 最\emph{小}的同样数量的路径；
  \item \textbf{不加路径}：只用问题和原始上下文，提示格式不变。
\end{itemize}
其中，“随机”和“反向”用来检验 $\Delta H$ 排序本身是否有价值；“不加路径”用来衡量整个知识图谱分支的净贡献。
我们还在 $\tau\in\{0,0.5,1,2,3\}$ 上扫描论文规则，统计“保留路径中含答案”和“一条路径都没保留”的题目比例。

\begin{figure}[htbp]
  \centering
  \begin{minipage}{0.48\textwidth}\centering
    \clfig{figure/cl_method/e2_accuracy.pdf}{\linewidth}
  \end{minipage}\hfill
  \begin{minipage}{0.48\textwidth}\centering
    \clfig{figure/cl_method/e2_tau_coverage.pdf}{\linewidth}
  \end{minipage}
  \caption{左：不同过滤策略的平均准确率；右：论文规则下阈值 $\tau$ 对保留结果的影响}
  \label{fig:cl-e2}
\end{figure}

\cltodo{表：各策略在各数据集上的准确率（来自 \texttt{e2\_accuracy.tex}）}

\cltodo{结果分析：代码规则与“随机”、“反向”、“不加路径”之间的差距；论文规则在 $\tau=1$ 时空集的比例；据此判断 H2}

% ---------------------------------------------------------------------
\subsubsection{流水线逐级损失与错误归因（E3）}
\label{sec:cl-e3}

为定位答案在哪一步丢失，我们逐题检查答案是否依次出现在：
S1 原始检索文本、S2 构建出的知识图谱、S3 候选路径、S4 被保留的路径中，以及 S5 最终是否答对。
由于最终提示中仍包含完整的原始上下文（第~\ref{sec:cl-modules}~小节），S5 并不受 S4 约束。
因此我们额外比较“保留路径含答案”与“保留路径不含答案”两种情况下的准确率，以区分模型的正确答案究竟来自路径，还是来自原文。

\begin{figure}[htbp]
  \centering
  \clfig{figure/cl_method/e3_funnel.pdf}{\textwidth}
  \caption{各数据集上答案的逐级保留率（虚线为标准 RAG 的准确率）}
  \label{fig:cl-e3}
\end{figure}

\cltodo{结果分析：最大的一级损失发生在哪里（S1→S2 抽取，还是 S2→S3 检索）；反事实数据上 S1→S2 的损失是否明显高于 golden 数据（即 A3 中的“抽取时被模型先验改正”）；S4 不成立时的准确率}

\cltodo{（P1）人工归因：从答错的样本中抽取 30 条，归入抽取错误、检索遗漏、过滤误删、生成错误四类，并给出一个典型案例}
          % X.4.1–X.4.2  Sean
% 补充实验 R3–R4—— 负责人：Hakimi
% （占位文件：由 Hakimi 撰写。标签 sec:cl-r3 / sec:cl-r4 已被其它小节引用，请保留。）
\subsubsection{外部知识错误下的鲁棒性}
\label{sec:cl-r3}
\cltodo{Hakimi：用 negative 数据的记录、以 gold\_true 重新评估；检验假设 H4（表~\ref{tab:cl-hypo}）}

\subsubsection{统计显著性与计算成本}
\label{sec:cl-r4}
\cltodo{Hakimi：McNemar 检验、bootstrap 置信区间、准确率与成本的关系}
      % X.4.3–X.4.4  Hakimi

\subsection{综合讨论：局限与改进方向}
\label{sec:cl-discussion}
% X.5 前半：方法层面的局限与改进 —— 负责人：Sean（D2 撰写）
\cltodo{Sean：结合 E1–E3 的结果，讨论方法层面的局限，并提出改进方向（例如：用“答案分布是否翻转”替代熵变化作为冲突信号；先在原文中定位冲突句、只抽取相关三元组以降低建图成本；对外部知识增加可信度校验以防御 A1 失效）}
     % X.5 前半  Sean
% 局限与展望：实验层面 —— 负责人：Hakimi
\subsubsection{实验层面}
\label{sec:cl-discussion-exp}
\cltodo{Hakimi：实验层面的局限（样本规模、单一模型、评测指标等）}
 % X.5 后半  Hakimi

% 小结 —— 负责人：Hakimi
\subsection{小结}
\label{sec:cl-summary}
\cltodo{Hakimi：半页小结}
                % X.6   Hakimi
